\section{Zariski's Main Theorem (algebraic version)}
\label{section-Zariski}

\noindent
This is the version you can prove using purely algebraic methods.
Before we can prove more powerful versions (for non-affine morphisms)
we need to develop more tools. See Cohomology of Schemes, Section
\ref{coherent-section-applications-formal-functions}
and
More on Morphisms, Section
\ref{more-morphisms-section-application-etale-neighbourhoods}.

\begin{theorem}[Algebraic version of Zariski's Main Theorem]
\label{theorem-main-theorem}
Let $f : Y \to X$ be an affine morphism of schemes.
Assume $f$ is of finite type.
Let $X'$ be the normalization of $X$ in $Y$. Picture:
$$
\xymatrix{
Y \ar[rd]_f \ar[rr]_{f'} & & X' \ar[ld]^\nu \\
& X &
}
$$
Then there exists an open subscheme $U' \subset X'$ such that
\begin{enumerate}
\item $(f')^{-1}(U') \to U'$ is an isomorphism, and
\item $(f')^{-1}(U') \subset Y$ is the set of points at which
$f$ is quasi-finite.
\end{enumerate}
\end{theorem}

\begin{proof}
There is an immediate reduction to the case where $X$ and hence $Y$
are affine. Say $X = \Spec(R)$ and $Y = \Spec(A)$.
Then $X' = \Spec(A')$, where $A'$ is the integral closure of
$R$ in $A$, see Definitions \ref{definition-integral-closure}
and \ref{definition-normalization-X-in-Y}. By
Algebra, Theorem \ref{algebra-theorem-main-theorem}
for every $y \in Y$ at which $f$ is quasi-finite, there exists an
open $U'_y \subset X'$ such that $(f')^{-1}(U'_y) \to U'_y$
is an isomorphism. Set $U' = \bigcup U'_y$ where $y \in Y$ ranges
over all points where $f$ is quasi-finite. It remains to show that
$f$ is quasi-finite at all points of $(f')^{-1}(U')$.
If $y \in (f')^{-1}(U')$ with image $x \in X$, then we see that
$Y_x \to X'_x$ is an isomorphism in a neighbourhood of $y$. Hence
there is no point of $Y_x$ which specializes to $y$, since this
is true for $f'(y)$ in $X'_x$, see Lemma \ref{lemma-integral-fibres}.
By Lemma \ref{lemma-quasi-finite-at-point-characterize} part (3)
this implies $f$ is quasi-finite at $y$.
\end{proof}

\noindent
We can use the algebraic version of Zariski's Main Theorem to show that
the set of points where a morphism is quasi-finite is open.

\begin{lemma}
\label{lemma-quasi-finite-points-open}
\begin{slogan}
The locally quasi-finite locus of a morphism is open
\end{slogan}
Let $f : X \to S$ be a morphism of schemes.
The set of points of $X$ where $f$ is quasi-finite is an open
$U \subset X$. The induced morphism $U \to S$ is locally quasi-finite.
\end{lemma}

\begin{proof}
Suppose $f$ is quasi-finite at $x$.
Let $x \in U = \Spec(A) \subset X$, $V = \Spec(R) \subset S$
be affine opens as in Definition \ref{definition-quasi-finite}.
By either Theorem \ref{theorem-main-theorem} above or
Algebra, Lemma \ref{algebra-lemma-quasi-finite-open},
the set of primes $\mathfrak q$ at which $R \to A$ is quasi-finite
is open in $\Spec(A)$. Since these all correspond to points
of $X$ where $f$ is quasi-finite we get the first statement.
The second statement is obvious.
\end{proof}

\noindent
We will improve the following lemma to general quasi-finite separated
morphisms later, see
More on Morphisms, Lemma
\ref{more-morphisms-lemma-quasi-finite-separated-pass-through-finite}.

\begin{lemma}
\label{lemma-quasi-finite-affine}
Let $f : Y \to X$ be a morphism of schemes.
Assume
\begin{enumerate}
\item $X$ and $Y$ are affine, and
\item $f$ is quasi-finite.
\end{enumerate}
Then there exists a diagram
$$
\xymatrix{
Y \ar[rd]_f \ar[rr]_j & & Z \ar[ld]^\pi \\
& X &
}
$$
with $Z$ affine, $\pi$ finite and $j$ an open immersion.
\end{lemma}

\begin{proof}
This is
Algebra, Lemma \ref{algebra-lemma-quasi-finite-open-integral-closure}
reformulated in the language of schemes.
\end{proof}

\begin{lemma}
\label{lemma-image-nowhere-dense-quasi-finite}
Let $f : Y \to X$ be a quasi-finite morphism of schemes.
Let $T \subset Y$ be a closed nowhere dense subset of $Y$.
Then $f(T) \subset X$ is a nowhere dense subset of $X$.
\end{lemma}

\begin{proof}
As in the proof of Lemma \ref{lemma-image-nowhere-dense-finite} this
reduces immediately to the case where the base $X$ is affine.
In this case $Y = \bigcup_{i = 1, \ldots, n} Y_i$ is a finite union
of affine opens (as $f$ is quasi-compact). Since each $T \cap Y_i$
is nowhere dense, and since a finite union of nowhere dense sets is
nowhere dense (see
Topology, Lemma \ref{topology-lemma-nowhere-dense}),
it suffices to prove that the image $f(T \cap Y_i)$ is nowhere dense in $X$.
This reduces us to the case where both $X$ and $Y$ are affine. At this point
we apply Lemma \ref{lemma-quasi-finite-affine} above to get a diagram
$$
\xymatrix{
Y \ar[rd]_f \ar[rr]_j & & Z \ar[ld]^\pi \\
& X &
}
$$
with $Z$ affine, $\pi$ finite and $j$ an open immersion.
Set $\overline{T} = \overline{j(T)} \subset Z$. By
Topology, Lemma \ref{topology-lemma-image-nowhere-dense-open}
we see $\overline{T}$ is nowhere dense in $Z$.
Since $f(T) \subset \pi(\overline{T})$
the lemma follows from the corresponding result in the finite case, see
Lemma \ref{lemma-image-nowhere-dense-finite}.
\end{proof}




